%
%  Description of VERARun
%
%
%%%%%%%%%%%%%%%%%%%%%%%%%%%%%%%%%%%%%%%%%
\section{VERARun}  \label{chap:verarun}

This chapter describes running cases with the VERARun script.
VERARun is the driver script that runs the input processor and corresponding VERA component codes.
VERARun also submits the job to the parallel job queue.

%%%%%%%%%%%%%%%%%%%%%%%%%%%%%%%%
\subsection{Running a Case}

VERARun is run by specifying {\tt verarun} on the command line, followed by the name of the input file.
There are additional command line options that are shown below.
\begin{Verbatim}[fontsize=\small]
   verarun <input file>
\end{Verbatim}

For example, if you input deck is called ``2a.inp'', you would enter:
\begin{Verbatim}[fontsize=\small]
   verarun 2a
\end{Verbatim}

Additional command line options can be found by typing {\tt verarun} with no arguments.
Doing so will return the following options.

\begin{Verbatim}[fontsize=\small]
usage: verarun [-x] [-e email_addr] [-h] [-c config_file] [-N job_name] [-l]
               [-n nprocs] [-O] [--ppn cpus_per_node] [-m mem_per_process]
               [-s subdir] [-d vera_install_dir] [-v vera_version] [--verbose]
               [-W] [-w job_id] [-t walltime] [--chain] [--debug]
               [--hostname host] [-r {overwrite,readwrite}]
               [--vera-installs-dir vera_installs_dir]
               [input_path [input_path ...]]

Creates and optionally submits machine-specific VERA jobs.

positional arguments:
  input_path            path to VERA input (.inp) or XML (.xml) files

optional arguments:
  -x, --dry-run         dry run only, create but don't execute the PBS script
  -e email_addr, --email-addrs email_addr
                        comma-delimited list of email addresses to notify of
                        job completion, defaults to ${USER}@$(hostname)
  -h, --help            print detailed help
  -c config_file, --host-config-file config_file
                        override host configuration file, supercedes
                        --hostname and --vera-installs-dir
  -N job_name, --job-name job_name
                        name for the PBS job
  -l, --list-vera-versions
                        list available VERA versions
  -n nprocs, --np nprocs, --nprocs nprocs, --num-procs nprocs
                        total number of processors need for the job (mpiexec
                        -np param), defaults to value computed from input
  -O, --output-job-name
                        print the job name to stdout
  --ppn cpus_per_node, --pnode cpus_per_node
                        specify processors per node, by default this is
                        calculated
  -m mem_per_process, --pmem mem_per_process, --proc-memory mem_per_process
                        specify memory required per processor in GB, defaults
                        to undefined
  -s subdir, --subdir subdir
                        create subdir, a value of "." specifies automatically
                        generated subdir name
  -d vera_install_dir, --vera-dir vera_install_dir
                        path to VERA installation, superceding --vera-
                        installs-dir, --vera-version, and the host
                        configuration
  -v vera_version, --vera-version vera_version
                        name of VERA version to use
  --verbose             turn on verbose messaging
  -W                    wait on job last submitted via verarun, overrides -W
  -w job_id, --wait-job-id job_id
                        ID of job which must complete before starting this job
  -t walltime, --wall-time walltime
                        wallclock execution time in floating point hours,
                        defaults to 24.0

advanced arguments:
  --chain, --chain-jobs
                        each case depends on its predecessor
  --debug               debug mode
  --hostname host       force the hostname
  -r {overwrite,readwrite}, --restart {overwrite,readwrite}
                        optional restart mode
  --vera-installs-dir vera_installs_dir
                        path to vera_installs directory containing VERA
                        versions

\end{Verbatim}

%%%%%%%%%%%%%%%%%%%%%%%%%%%%%%%%
\subsection{VERARun Output}

Upon completion of a VERA job, several output files may be created depending on the code options used.
Some typical outputs include:
\begin{itemize}
\item VERAIn XML file.  This file is written upon the successful completion of VERAIn
\item VERA HDF output file.  This is a binary file with results that can be visualized in VERAView, or post-processed with user utility codes.
\item MPACT output file.  This file is written upon the successful completion of MPACT (if applicable).
\item MPACT log File.  This file is written upon the successful completion of MPACT (if applicable).
\item MPACT summary File.  This file is written upon the successful completion of MPACT (if applicable).
\item Standard output file. This file is a log of all output written to the standard output.
\item Standard error file.  this file is a log of all output written to the standard error file.
\end{itemize}



%%%%%%%%%%%%%%%%%%%%%%%%%%%%%%%%
\subsection{Input Errors}

If the {\tt verarun} command does not work, the user should make sure that it is in the path.
The user may need to consult with the system administrator for the correct path.

The next step when looking for an error is to look at the standard error file.
If the job ran successfully, the size of this file will be zero.

If there were any errors in the input processor (VERAIn), the errors will be written to the standard error file.
Common errors from the input processor include
\begin{enumerate}
\item invalid keywords
\item invalid map sizes
\item invalid input options
\end{enumerate}

If the input processor works correctly, an error may still occur in one of the VERA component codes.
The user used look at the error message and consult the user manual for the component code.












